%!TEX program = xelatex
\def\numeromodulo{2}

\def\numerolista{4}

\documentclass[12pt]{exam}

\def\ano{2023}
\def\semestre{1}
\def\disciplina{Introdução à \'Algebra Linear}
\def\turma{3}


\usepackage{caption}
\usepackage{amssymb}
\usepackage{amsmath,amsfonts,amsthm,amstext}
\usepackage[utf8]{inputenc} % Required for inputting international characters
\usepackage[T1]{fontenc}
\usepackage[brazil]{babel}
\usepackage{graphicx}
\graphicspath{{../figuras_geogebra/}{../Pictures/}}
\usepackage{enumitem}
\usepackage{multicol}
\usepackage{answers}
\usepackage[svgnames]{xcolor}
\usepackage{tikz}
\usepackage{ifthen}
\usetikzlibrary{lindenmayersystems}
\usetikzlibrary[shadings]
\usepackage{tikz-3dplot}
\usepackage{tcolorbox}
%-------------amatrix
% Augmented matrix.  Usage (note the argument does not count the aug col):
% \begin{amatrix}{2}
%   1  2  3 \\  4  5  6
% \end{amatrix}
\newenvironment{amatrix}[1]{%
  \left(\begin{array}{@{}*{#1}{r}|r@{}}
}{%
  \end{array}\right)
}


\Newassociation{solucao}{Solution}{ans}
\newtheorem{exercicio}{}

\setlength{\topmargin}{-1.0in}
\setlength{\oddsidemargin}{0in}
\setlength{\textheight}{10.1in}
\setlength{\textwidth}{6.5in}
\setlength{\baselineskip}{12mm}

\newcounter{exercicios}
\setcounter{exercicios}{0}
\newcommand{\questao}{
\addtocounter{exercicios}{1}
\noindent{\bf Exerc{\'\i}cio \Roman{exercicios}: }}

\newcommand{\resp}[1]{
\noindent{\bf Exerc{\'\i}cio #1: }}

\extraheadheight{0.7in}
\firstpageheadrule
\runningheadrule
\lhead{
        \begin{minipage}[c]{1.7cm}
        \includegraphics[width=1.7cm]{unb.pdf}
        \end{minipage}%
        \hspace{0pt}
        \begin{minipage}[c]{4in}
          {Universidade de Bras{\'\i}lia} --
          {Departamento de Matem{\'a}tica}
\end{minipage}
\vspace*{-0.8cm}
}
% \chead{Universidade de Bras{\'\i}lia - Departamento de Matem{\'a}tica}
% \rhead{}
% \vspace*{-2cm}

\extrafootheight{.5in}
\footrule
\lfoot{\disciplina\ - \semestre$^o$/\ano\ - Módulo \numeromodulo}
\cfoot{}
\rfoot{P\'agina \thepage\ de \numpages}

\newcommand{\nsub}{\varsubsetneq}
\newcommand{\vaz}{\emptyset}
\newcommand{\im}{{\rm Im\,}}
\newcommand{\sub}{\subseteq}
\renewcommand{\sin}{{\rm sen\,}}
\newcommand{\n}{\mathbb{N}}
\newcommand{\z}{\mathbb{Z}}
\newcommand{\rac}{\mathbb{Q}}
\newcommand{\real}{\mathbb{R}}
\newcommand{\complex}{\mathbb{C}}
\newcommand{\cp}[1]{\mathbb{#1}}
\newcommand{\ch}{\mbox{\textrm{car\,}}\nobreak}
\newcommand{\vesp}[1]{\vspace{ #1  cm}}
\newcommand{\compcent}[1]{\vcenter{\hbox{$#1\circ$}}}
\newcommand{\comp}{\mathbin{\mathchoice
{\compcent\scriptstyle}{\compcent\scriptstyle}
{\compcent\scriptscriptstyle}{\compcent\scriptscriptstyle}}}


\newcommand{\integer}{\mathbb{Z}}

\begin{document}

    \Opensolutionfile{ans}[ans1]

    \begin{center}
        {\Large\bf \disciplina\ - Turma \turma\ -- \semestre$^{o}$/\ano} \\ \vspace{9pt} {\large\bf
            $\numerolista^a$ Lista de Exercícios -- Módulo \numeromodulo}\\ \vspace{9pt} Prof. José Antônio O. Freitas
    \end{center}

    \hrule

    \begin{exercicio}
        Nos itens abaixo, encontre o vetor de coordenadas de $w$ em relação à base  $S = \{v_1, v_2\}$ de $\real^2$:
        \begin{enumerate}[label={\alph*})]
            \item $v_1 = (1, 0)$, $v_2 = (0, 2)$; $w = (3, -7)$
            \item $v_1 = (2, -4)$, $v_2 = (3, 8)$; $w = (1, 1)$
            \item $v_1 = (1, 1)$, $v_2 = (0, 2)$; $w = (a, b)$
        \end{enumerate}
    \end{exercicio}

    \begin{exercicio}
        Nos itens abaixo, encontre o vetor de coordenadas de $w$ em relação à base $S = \{v_1, v_2\}$ de $\real^2$:
        \begin{enumerate}[label={\alph*})]
            \item $v_1 = (1, -1)$, $v_2 = (1, 1)$; $w = (1, 0)$
            \item $v_1 = (1, -1)$, $v_2 = (1, 1)$; $w = (0, 1)$
            \item $v_1 = (1, -1)$, $v_2 = (1, 1)$; $w = (1, 1)$
        \end{enumerate}
    \end{exercicio}

    \begin{exercicio}
        Nos itens abaixo, encontre o vetor de coordenadas de $w$ em relação à base $S = \{v_1, v_2, v_3\}$ de $\real^3$:
        \begin{enumerate}[label={\alph*})]
            \item $v_1 = (1, 0, 0)$, $v_2 = (2, 2, 0)$, $v_3 = (3, 3, 3)$; $w = (2, -1, 3)$
            \item $v_1 = (1, 2, 3)$, $v_2 = (-4, 5, 6)$, $v_3 = (7, -8, 9)$; $w = (5, -12, 3)$
        \end{enumerate}
    \end{exercicio}

    \begin{exercicio}
        Nos itens abaixo, encontre o vetor de coordenadas de $q(x)$ em relação à base $S = \{p_1(x), p_2(x), p_3(x)\}$ de $\mathcal{P}_2(\real)$:
        \begin{enumerate}[label={\alph*})]
            \item $p_1(x) = 1$, $p_2(x) = x$, $p_3(x) = x^2$; $q(x) = 4 - 3x + x^2$
            \item $p_1(x) = 1 + x$, $p_2(x) = 1 + x^2$, $p_3(x) = x + x^2$; $q(x) = 2 - x + x^2$
        \end{enumerate}
    \end{exercicio}

    \begin{exercicio}
        Encontre o vetor de coordenadas da matriz $A$ em relação à base $S = \{A_1, A_2, A_3, A_4\}$ de $M_2(\real)$:
        \begin{align*}
            A = \begin{bmatrix}
                \phantom{-} 2 & 0\\
                -1 & 3
            \end{bmatrix},
            A_1 = \begin{bmatrix}
                -1 & 1\\
                \phantom{-} 0 & 0
            \end{bmatrix},
            A_2 = \begin{bmatrix}
                1 & 1\\
                0 & 0
            \end{bmatrix},
            A_3 = \begin{bmatrix}
                0 & 0\\
                1 & 0
            \end{bmatrix},
            A_4 = \begin{bmatrix}
                0 & 0\\
                0 & 1
            \end{bmatrix}.
        \end{align*}
    \end{exercicio}


    \begin{exercicio}
        Encontre o vetor de coordenadas de
        \[
            A = \begin{bmatrix}
                1 & 2\\
                3 & 4
            \end{bmatrix}
        \]
        em relação à base $\mathcal{B} = \{E_{22}, E_{21}, E_{12}, E_{11}\}$ de $M_2(\real)$.
    \end{exercicio}

    \begin{exercicio}
        Encontre o vetor de coordenadas da matriz $A$ em relação à base $S = \{A_1, A_2, A_3, A_4\}$ de $M_2(\real)$:
        \begin{align*}
            A = \begin{bmatrix}
                1 & 2\\
                3 & 4
            \end{bmatrix},
            A_1 = \begin{bmatrix}
                1 & 0\\
                0 & 0
            \end{bmatrix},
            A_2 = \begin{bmatrix}
                1 & 1\\
                0 & 0
            \end{bmatrix},
            A_3 = \begin{bmatrix}
                1 & 1\\
                1 & 0
            \end{bmatrix},
            A_4 = \begin{bmatrix}
                1 & 1\\
                1 & 1
            \end{bmatrix}.
        \end{align*}
    \end{exercicio}

    \begin{exercicio}
        Encontre o vetor de coordenadas de $p(x) = 1 + 2x + 3x^2$ em relação à base $S = \{1 + x, 1 - x, x^2\}$ de $\mathcal{P}_2(\real)$.
    \end{exercicio}

    \begin{exercicio}
        Encontre o vetor de coordenadas de $p(x) = 2 - x + 3x^2$ em relação à base $S = \{1, 1 + x, -1 + x^2\}$ de $\mathcal{P}_2(\real)$.
    \end{exercicio}

    \begin{exercicio}
        Encontre uma base de $\real^4$ contendo os vetores $u = (1, -4, 2, -3)$ e $v = (-3, 8, -4, 6)$.
    \end{exercicio}

    \begin{exercicio}
        Os vetores $v_1 = (1, -2, 3, -5)$ e $v_2 = (0, -1, 2, -3)$ são linearmente independentes. Estenda $\{v_1, v_2\}$ até uma base de $\real^4$.
    \end{exercicio}

    \begin{exercicio}
        Estenda o conjunto $\{1 + x, 1 + x^2 + x^3\}$ a uma base de $\mathcal{P}_3(\real)$.
    \end{exercicio}

    \begin{exercicio}
        Estenda o conjunto
        \[
        S = \left\{
            \begin{bmatrix}
                0 & 1 & \phantom{-} 1\\
                1 & 0 & -1\\
                0 & 0 & \phantom{-} 2
            \end{bmatrix},
            \begin{bmatrix}
                -1 & 0 & \phantom{-} 1\\
                \phantom{-} 0 & 2 & -1\\
                \phantom{-} 0 & 1 & \phantom{-} 2
            \end{bmatrix}
        \right\}
        \]
        a uma base de $M_3(\real)$.
    \end{exercicio}

    \begin{exercicio}
        Encontre uma base para $Span(1, 1 + x, 2x)$ em $\mathcal{P}_2(\real)$.
    \end{exercicio}

    \begin{exercicio}
        Encontre uma base para $Span(1 - 2x, 2x - x^2, 1 - x^2, 1 + x^2)$ em $\mathcal{P}_2(\real)$.
    \end{exercicio}

    \begin{exercicio}
        Encontre uma base para $Span(1 - x, x - x^2, 1 - x^2, 1 - 2x + x^2)$ em $\mathcal{P}_2(\real)$.
    \end{exercicio}

    \begin{exercicio}
        Encontre uma base para $Span(A_1, A_2, A_3, A_4)$ em $M_2(\real)$, onde:
        \[
            A_1 = \begin{bmatrix}
                1 & 0\\
                0 & 1
            \end{bmatrix},
            A_2 = \begin{bmatrix}
                0 & 1\\
                1 & 0
            \end{bmatrix},
            A_3 = \begin{bmatrix}
                -1 & \phantom{-} 1\\
                \phantom{-} 1 & -1
            \end{bmatrix},
            A_4 = \begin{bmatrix}
                \phantom{-} 1 & -1\\
                -1 & \phantom{-} 1
            \end{bmatrix}.
        \]
    \end{exercicio}


    %\newpage

    %\Closesolutionfile{ans}

    %\begin{center}
    %    {\large\bf RESPOSTAS}
    %\end{center}

    %\hrule

    %\input{ans1}

\end{document}
